\chapter{Yellow Tongue}

\section*{Definition}
Yellowish discoloration of the dorsal tongue surface, usually resulting from accumulation of bacteria, debris, or staining on elongated filiform papillae.

\section*{Pathophysiology}
Yellow tongue typically arises from overgrowth of chromogenic bacteria or fungi (often Candida) on the tongue dorsum, especially when filiform papillae are elongated (hairy tongue). Elongation occurs when normal desquamation is impaired by poor oral hygiene, dry mouth, soft diet, or medications. The yellow color may also result from staining by food, tobacco, or mouthwashes containing oxidizing agents.

\section*{Etiology}
\begin{itemize}
  \item Poor oral hygiene (bacterial overgrowth)
  \item Dry mouth (xerostomia)
  \item Dehydration
  \item Smoking or tobacco use
  \item Mouthwashes containing peroxide or witch hazel
  \item Oral thrush (Candida overgrowth, can appear yellow)
  \item Coffee, tea, or food staining
  \item Antibiotic use (altered oral flora)
  \item Fever or illness (reduced oral clearance)
  \item Jaundice (rare — sclera also yellow, bilirubin elevated)
\end{itemize}

\section*{Indications for Evaluation}
Seek care if:
\begin{itemize}
  \item Severe or worsening symptoms
  \item Yellow tongue persisting more than 2 weeks despite good oral hygiene, or with jaundice, pain, or difficulty eating
  \item Accompanied by other concerning signs
\end{itemize}

\section*{Related Symptoms}
\begin{itemize}
  \item Bad breath
  \item White tongue
  \item Dry mouth
  \item Taste changes
  \item Oral thrush
\end{itemize}

\textbf{Note:} Always consider serious underlying conditions when evaluating persistent yellow tongue.

\section*{Self-Care / Home Management}
\begin{itemize}
  \item Rest and avoid activities that may aggravate the condition.
  \item Maintain adequate hydration and a balanced diet.
  \item Over-the-counter remedies may provide symptomatic relief (consult a pharmacist or doctor).
  \item Monitor symptoms and seek professional advice if they persist or worsen.
  \item Avoid self-medicating with prescription medications without medical guidance.
\end{itemize}

\section*{Diagnostic Approach}
\begin{itemize}
  \item A thorough history and physical examination are the first steps in evaluation.
  \item Laboratory tests (blood work, urinalysis) may be ordered based on clinical suspicion.
  \item Imaging studies (X-ray, ultrasound, CT, MRI) may be indicated when structural pathology is suspected.
  \item Specialist referral is arranged when the diagnosis remains uncertain or requires specific expertise.
\end{itemize}

\section*{Treatment / Management Overview}
\begin{itemize}
  \item Treatment targets the underlying cause identified through diagnostic evaluation.
  \item Supportive care and symptom management are provided while awaiting definitive therapy.
  \item Medications, lifestyle modifications, and physical therapies may be prescribed as appropriate.
  \item Follow-up ensures treatment effectiveness and allows timely adjustments to the care plan.
\end{itemize}

\clearpage